\documentclass[11pt]{article}
\usepackage[margin=1in]{geometry}
\usepackage{enumitem}
% =============================================================================
% Preambles/header.tex — shared LaTeX/Beamer preamble for this project
%
% Use from a Beamer lecture by prepending:
%     \input{header}
% and compiling with TEXINPUTS=../Preambles:$TEXINPUTS (the /compile-latex
% skill does this automatically).
%
% PALETTE CONTRACT. Color names defined here MUST match the names in
% Quarto/theme-template.scss so Beamer and Quarto renderings use the same
% palette. The scripts/check-palette-sync.sh script enforces this.
%
% To customize: edit the HEX values below AND the matching $variable in
% Quarto/theme-template.scss, then run scripts/check-palette-sync.sh.
% =============================================================================

% -----------------------------------------------------------------------------
% Engine detection + encoding
%
% Our /compile-latex skill uses XeLaTeX, where inputenc is unnecessary and
% can produce encoding warnings. Only load inputenc under pdfTeX.
% -----------------------------------------------------------------------------
\usepackage{iftex}
\ifPDFTeX
  \usepackage[utf8]{inputenc}
\fi

\usepackage{amsmath, amssymb, amsthm}
\usepackage{booktabs}
\usepackage{graphicx}

% -----------------------------------------------------------------------------
% PALETTE — mirrors Quarto/theme-template.scss variable names.
% Load xcolor and define colors BEFORE anything that references them
% (notably hyperref, hypersetup, and the Beamer color assignments).
% -----------------------------------------------------------------------------
\usepackage{xcolor}

\definecolor{primary-blue}{HTML}{012169}      % headings, accents
\definecolor{primary-gold}{HTML}{B9975B}      % emphasis, borders
\definecolor{highlight-yellow}{HTML}{F2A900}  % markers, alerts
\definecolor{light-bg}{HTML}{E8EDF5}          % title / section backgrounds
\definecolor{jet}{HTML}{1A1A1A}               % body text

% Semantic colors (used in diagrams and annotations)
\definecolor{positive}{HTML}{15803D}          % good / observed / identified
\definecolor{negative}{HTML}{B91C1C}          % bad / problematic / bias
\definecolor{neutral}{HTML}{525252}           % reference / context

% Highlight accents (match .hi-* classes in the SCSS)
\definecolor{hi-slate}{HTML}{314F4F}
\definecolor{hi-green}{HTML}{2E8B57}
\definecolor{hi-red}{HTML}{C41E3A}

% Semantic aliases so skill templates and snippets can write
% \textcolor{accent}{...} without hard-coding "primary-blue".
\colorlet{accent}{primary-blue}
\colorlet{accent2}{primary-gold}

% -----------------------------------------------------------------------------
% Hyperref — loaded AFTER xcolor + palette so link colors resolve.
% -----------------------------------------------------------------------------
\usepackage{hyperref}
\hypersetup{colorlinks=true, linkcolor=primary-blue, urlcolor=primary-blue,
            citecolor=primary-blue, pdfborder={0 0 0}}

% -----------------------------------------------------------------------------
% Course code — set once per deck (after \input{header}, before \begin{document})
% via \coursecode{CS401}. Shown in the footer on every slide so a deck's course
% is identifiable without opening the title page. Empty by default.
% -----------------------------------------------------------------------------
\makeatletter
\newcommand{\coursecode}[1]{\gdef\@coursecodetext{#1}}
\gdef\@coursecodetext{}
\makeatother

% -----------------------------------------------------------------------------
% Beamer theme assignments (only apply under Beamer).
%
% We detect Beamer via \@ifclassloaded{beamer}, which is the documented,
% non-internal way. This must be wrapped in \makeatletter/\makeatother
% because \@ifclassloaded contains an @-catcode character.
% -----------------------------------------------------------------------------
\makeatletter
\@ifclassloaded{beamer}{%
  \usefonttheme{professionalfonts}
  \setbeamercolor{structure}{fg=primary-blue}
  \setbeamercolor{frametitle}{fg=primary-blue}
  \setbeamercolor{title}{fg=primary-blue}
  \setbeamercolor{subtitle}{fg=primary-gold}
  \setbeamercolor{author}{fg=primary-blue}
  \setbeamercolor{institute}{fg=neutral}
  \setbeamercolor{date}{fg=primary-gold}
  \setbeamercolor{itemize item}{fg=primary-blue}
  \setbeamercolor{itemize subitem}{fg=primary-gold}
  \setbeamercolor{alerted text}{fg=primary-gold}
  \setbeamercolor{block title}{fg=white, bg=primary-blue}
  \setbeamercolor{block body}{bg=light-bg}
  % Semantic block variants: block=definition/concept (blue), exampleblock=
  % worked example (green/positive), alertblock=key takeaway or caution (gold).
  % Keep to <=2 colored boxes per slide across ALL three variants (INV-7).
  \setbeamercolor{block title example}{fg=white, bg=positive}
  \setbeamercolor{block body example}{bg=light-bg}
  \setbeamercolor{block title alerted}{fg=white, bg=primary-gold}
  \setbeamercolor{block body alerted}{bg=light-bg}

  % Minimal footer: course code (left) + page X / N (right), no navigation clutter
  \setbeamertemplate{navigation symbols}{}
  \setbeamertemplate{footline}{%
    \color{neutral}\footnotesize\hspace{0.7em}\@coursecodetext\hfill
    \insertframenumber/\inserttotalframenumber\hspace{0.7em}\vspace{0.3em}}
}{}
\makeatother

% -----------------------------------------------------------------------------
% TikZ setup — libraries the snippet gallery uses
% -----------------------------------------------------------------------------
\usepackage{tikz}
\usetikzlibrary{arrows.meta, positioning, calc, decorations.pathreplacing,
                fit, shapes.geometric, backgrounds}

% Shared TikZ styles — snippets and hand-written diagrams can pick these up
% via `\begin{tikzpicture}[every node/.style={...}]` or by naming specific
% styles (e.g., `\node[dag-node] (X) at (0,0) {X};`).
\tikzset{
  % Node styles for causal DAGs and similar
  dag-node/.style = {
    draw=primary-blue, fill=light-bg, rounded corners=2pt,
    minimum width=1.2cm, minimum height=0.9cm, align=center, font=\small
  },
  decision-node/.style = {
    draw=primary-gold, fill=white, diamond, aspect=1.8,
    minimum width=1.8cm, minimum height=1.2cm, align=center, font=\small,
    inner sep=1pt
  },
  % Edge styles — observed vs counterfactual
  observed-edge/.style = {-{Latex[length=2.5mm]}, thick, draw=jet},
  counterfactual-edge/.style = {-{Latex[length=2.5mm]}, thick, draw=neutral, dashed},
  confound-edge/.style = {-{Latex[length=2.5mm]}, thick, draw=negative, dashed},
  % Data-point styles for plots
  observed-dot/.style = {fill=primary-blue, draw=primary-blue, circle, inner sep=1.2pt},
  counterfactual-dot/.style = {fill=white, draw=neutral, circle, inner sep=1.2pt}
}

% -----------------------------------------------------------------------------
% Convenience macros
% -----------------------------------------------------------------------------
\newcommand{\muted}[1]{\textcolor{neutral}{#1}}
\newcommand{\key}[1]{\textcolor{primary-gold}{\textbf{#1}}}
\newcommand{\good}[1]{\textcolor{positive}{#1}}
\newcommand{\bad}[1]{\textcolor{negative}{#1}}

% Standout frame macro: full-bleed dark background for section transitions.
% Beamer-only (uses \begin{frame}/\setbeamercolor/\usebeamerfont) -- do not
% call from an article-class file (e.g. Notes/<CODE>/*-notes.tex). \muted,
% \key, \good, \bad, and \coursecode above are plain xcolor/gdef and are
% safe to use from either class; verified when Notes/article-class support
% was added.
% Expand: \transitionslide{Your transition title}
\newcommand{\transitionslide}[1]{%
  \begingroup
  \setbeamercolor{background canvas}{bg=primary-blue}
  \begin{frame}[plain]
    \centering
    \vfill
    {\usebeamerfont{title}\color{white}\Large #1\par}
    \vfill
  \end{frame}
  \endgroup
}

% Section divider frame (Beamer-only), an alternative to \transitionslide with a
% darker two-line style: near-black (jet) background, a small white label line
% above a larger gold title, and a thin gold rule along the bottom edge.
% Expand: \sectiondivider{Section 2}{Pointers}
\newcommand{\sectiondivider}[2]{%
  \begingroup
  \setbeamercolor{background canvas}{bg=jet}
  \begin{frame}[plain]
    \vspace*{3.0cm}
    {\color{white}\ttfamily\large #1\par}
    \vspace{0.5cm}
    {\color{primary-gold}\LARGE #2\par}
    \begin{tikzpicture}[remember picture, overlay]
      \draw[primary-gold, line width=2.5pt]
        ($(current page.south west)+(0,0.1)$) -- ($(current page.south east)+(0,0.1)$);
    \end{tikzpicture}
  \end{frame}
  \endgroup
}

% -----------------------------------------------------------------------------
% Channel watermark — "Concepts That Click" (YouTube).
%
% Article-class files (CompetitiveExam Q/A sets, Notes, Assignments) get a
% light diagonal draft watermark on every page plus a centered footer naming
% the channel. Beamer decks get a subtle TikZ background watermark instead
% (the footline already carries the course code). Centralized here so every
% MCQ PDF is branded without per-file edits.
% -----------------------------------------------------------------------------
\makeatletter
\@ifclassloaded{article}{%
  \usepackage{draftwatermark}
  \SetWatermarkText{Concepts That Click}
  \SetWatermarkScale{1}
  \SetWatermarkFontSize{2cm}
  \SetWatermarkLightness{0.86}
  \SetWatermarkAngle{45}
  \usepackage{fancyhdr}
  \pagestyle{fancy}
  \fancyhf{}
  \fancyfoot[C]{\footnotesize\textcolor{neutral}{Concepts That Click~|~YouTube: \href{https://www.youtube.com/@concepts.thatclick}{@concepts.thatclick}}}
  \renewcommand{\headrulewidth}{0pt}
  \renewcommand{\footrulewidth}{0pt}
  \fancypagestyle{plain}{%
    \fancyhf{}%
    \fancyfoot[C]{\footnotesize\textcolor{neutral}{Concepts That Click~|~YouTube: \href{https://www.youtube.com/@concepts.thatclick}{@concepts.thatclick}}}%
    \renewcommand{\headrulewidth}{0pt}%
    \renewcommand{\footrulewidth}{0pt}%
  }
}{}
\@ifclassloaded{beamer}{%
  \setbeamertemplate{background}{%
    \begin{tikzpicture}[remember picture, overlay]
      \node[rotate=45, opacity=0.055, align=center]
        at (current page.center)
        {\Huge\textcolor{neutral}{Concepts That Click}};
    \end{tikzpicture}%
  }
}{}
\makeatother 
\coursecode{JKSSB}
\title{Lesson 15 Practice: Answers and Rationales}
\author{JKSSB Computer Awareness}
\date{\today}
\begin{document}
\maketitle
\noindent\textbf{Provenance:} All 20 items are original JKSSB-pattern
questions (\textit{[Original]}); none reproduces an official paper item.
\begin{enumerate}[label=\textbf{Q\arabic*.}, leftmargin=*]
\item \textbf{C.} A dot matrix printer strikes an inked ribbon against the
  paper, so it is an impact printer. Inkjet (A), laser (B), and thermal (D)
  form the image without striking, so they are non-impact.
\item \textbf{B.} Striking an inked ribbon against the paper is exactly the
  definition of an impact printer. A non-impact printer (A) does not strike
  the paper; a page printer (C) is classified by the unit printed, not by
  contact; a plotter (D) draws lines and is not defined this way.
\item \textbf{D.} The inkjet printer is non-impact because it sprays ink
  without touching the paper. The dot matrix (A), daisy wheel (B), and drum
  printer (C) all strike a ribbon or paper and are impact printers.
\item \textbf{B.} An inkjet printer sprays droplets of liquid ink through
  nozzles. A laser printer (A) uses toner; a thermal printer (C) uses heat;
  a dot matrix printer (D) strikes a ribbon.
\item \textbf{B.} A laser printer uses a laser beam and toner in a
  xerographic process. An inkjet printer (A) sprays liquid ink; a dot matrix
  printer (C) strikes a ribbon; a thermal printer (D) uses heat on thermal
  paper.
\item \textbf{C.} A laser printer prints a whole page at a time, so it is a
  page printer. A dot matrix (A) and a daisy wheel printer (B) are character
  printers; a plotter (D) draws lines and is not a page printer.
\item \textbf{C.} A thermal printer applies heat to specially treated paper,
  which is why it is used for receipts and labels. A laser printer (A) uses
  a laser beam and toner; an inkjet printer (B) sprays liquid ink; a dot
  matrix printer (D) strikes a ribbon.
\item \textbf{B.} A 3D printer works by additive manufacturing, adding
  material layer by layer to build an object. Striking a ribbon (A) is dot
  matrix, spraying ink (C) is inkjet, and melting toner (D) is laser --- all
  two-dimensional paper printing.
\item \textbf{B.} A plotter draws vector line graphics with pens, so it suits
  large engineering drawings, maps, and CAD plans. Inkjet (A) and laser (C)
  printers print dots/characters rather than drawn lines; a dot matrix printer
  (D) is a low-quality character printer.
\item \textbf{B.} An inkjet printer sprays liquid ink, while a laser printer
  uses a laser beam and toner --- the correct contrast. A reverses the
  consumables; C is wrong because both are non-impact; D swaps their print
  units (a laser is a page printer, not a character printer).
\item \textbf{B.} A printer that prints one whole line at a time, such as a
  drum or chain printer, is a line printer. A character printer (A) prints one
  character at a time; a page printer (C) prints a whole page; a plotter (D)
  draws lines and is not a printer of this kind.
\item \textbf{A.} DPI stands for Dots Per Inch, a measure of print/display
  sharpness. B, C, and D are not the correct expansions.
\item \textbf{B.} Resolution is the number of pixels a display shows,
  written as columns $\times$ rows. Refresh rate (A) is the redraws per
  second; response time (C) is the pixel switching time; contrast ratio (D)
  compares white and black levels.
\item \textbf{A.} A pixel (picture element) is the smallest addressable
  element of a display. A byte (B) is a unit of data; a bitmap (C) is an
  image representation; a raster (D) is a pixel grid, not a single element.
\item \textbf{A.} Refresh rate is measured in hertz (Hz), the number of
  redraws per second. Milliseconds (B) measure response time; pixels (C)
  measure resolution; nits (D) measure brightness.
\item \textbf{B.} Response time is measured in milliseconds (ms), the time a
  pixel takes to change state. Hertz (A) measures refresh rate; DPI (C)
  measures sharpness; nits (D) measure brightness.
\item \textbf{B.} The ratio between the brightest white and darkest black is
  the contrast ratio. Refresh rate (A) is redraws per second; resolution (C)
  is the pixel count; response time (D) is the pixel switching time.
\item \textbf{B.} Brightness, the light a display emits, is commonly measured
  in nits (cd/m$^2$). Hertz (A) measures refresh rate; milliseconds (C)
  measure response time; DPI (D) measures sharpness.
\item \textbf{D.} The dot matrix printer is an \textbf{impact} printer,
  because it strikes the ribbon; calling it non-impact is the incorrect
  statement. Inkjet (A), laser (B), and thermal (C) printers are genuinely
  non-impact.
\item \textbf{B.} Statements I and II are correct: a dot matrix printer is
  impact, and a laser printer is non-impact. III is false because a thermal
  printer is non-impact, not impact. A omits II; C includes false III and
  omits I; D includes false III.
\end{enumerate}
\end{document}
